\documentclass[12pt,a3paper]{article}
\usepackage{better_poster}
\usepackage{qrcode}

% Marks text that is still missing. Search for "FILL IN" before printing.
\newcommand{\fillin}[1]{\textcolor{red}{[FILL IN: #1]}}

% authors
\title{A tracked desk robot that shows emotions while it cleans your table}
\author{A project in IN5590 by Markus Yuki Nishimura and Martin Haugen\\
E-mail: \mbox{markusyn@uio.no}, \mbox{martin27@uio.no}}

% type of poster: [exp]erimental results, [methods], [theory]
\newcommand\postertype{exp}

\begin{document}

% header: the name of the robot, with the full name below in smaller text
\makefinding{
    \textbf{M.A.T.S}\\[0.3em]
    {\fontsize{24}{30}\selectfont Mechanized Assistant for Tabletop Sanitation}
}

\makemain{
    % space between paragraphs (the column box sets it to zero)
    \setlength{\parskip}{0.5em}
    \raggedcolumns
    \begin{multicols}{2}
        \section{Intro}
        Desks collect crumbs, dust and eraser debris. M.A.T.S drives on tracks and
        pushes the debris ahead of a brush to the table edge. There it lowers a bucket
        on two arms over the edge, and the brush pushes the debris into the bucket. A
        display in its head shows seven states: neutral, working, happy, sad,
        surprised, scared and angry.

        Our goals:
        \begin{itemize}
        \setlength\itemsep{0.1em}\setlength\parskip{0pt}
            \item \textbf{Physical design:} the bucket reaches past the table edge and
            stays upright, and the robot never tips.
            \item \textbf{Safety:} stop before a track leaves the table, lock the
            tracks while the bucket is lowered, and cut all motor power with a switch.
            \item \textbf{Mobility:} drive 1\,m forward and back, and turn
            360$^\circ$ in place.
            \item \textbf{Expression:} show the seven states as animations.
        \end{itemize}

        \section{Tools}
        We designed the robot in SolidWorks and printed the PLA parts on a Creality
        Ender-3 V3 SE. The lab engineer printed the TPU tracks. Key components:
        \begin{itemize}
        \setlength\itemsep{0.1em}\setlength\parskip{0pt}
            \item Arduino Nano ESP32 microcontroller
            \item 2 DAGU DG01D gearmotors with encoders
            \item 2 L293D motor driver chips
            \item 2 micro servos and a linear actuator for the arms
            \item 2 VL53L4CD time-of-flight edge sensors
            \item 2.8-inch ILI9341 display and an IMU
        \end{itemize}

    % this determines where the columns are separated
    \columnbreak

        \textbf{Software.} The firmware is written in C++17 on the Arduino-ESP32 core
        and built with PlatformIO. One core runs the low-level control loop: the edge
        reflex, the motor, servo and actuator control, and the dump interlock. The
        other core runs the cleaning behavior, the emotion state and the face. A track
        may move forward only while the edge sensor ahead of it sees the table, and the
        cleaning behavior can never override this. The display uses the Adafruit
        ILI9341 and GFX libraries, and the edge sensors use the STM32duino VL53L4CD
        library.

        \textbf{Optimization.} \fillin{what we optimized and how we did it}

        \section{Conclusion}
        \fillin{how far we came in reaching the goals}

        \fillin{the results of the optimization}

        \textbf{Future work.} \fillin{what would improve the robot, and what we would
        do differently}

    \end{multicols}
}
% figures (no blank line above, or this column moves below the main column)
\makeextracolumn{
    \includegraphics[width=\linewidth]{1assembled.PNG}

    {\small \textbf{Figure 1.} M.A.T.S assembled in SolidWorks.}

    \vspace{1em}
    \includegraphics[width=\linewidth]{2physical_prototype.jpg}

    {\small \textbf{Figure 2.} The printed prototype.}
}

% footer: UiO logo on the left, QR code to the project 2 repo on the right
\makealtfooter{images/uni_logo.pdf}{%
    \qrcode[height=0.11\paperheight]{https://github.uio.no/martin27/in5590\_project2}%
}

\end{document}
